% Configuração visual — TCE-MA Auditor TI 2026
\usepackage[T1]{fontenc}
\usepackage[utf8]{inputenc}
\usepackage[brazil]{babel}
\usepackage{lmodern}
\usepackage{microtype}
\usepackage{geometry}
\geometry{a4paper,top=2.2cm,bottom=2.2cm,left=2.1cm,right=2.1cm,headheight=15pt}
\usepackage{amsmath,amssymb,mathtools}
\usepackage{booktabs,longtable,tabularx,array,multirow}
\usepackage{graphicx}
\usepackage{enumitem}
\usepackage{xcolor}
\usepackage{tikz}
\usetikzlibrary{arrows.meta,positioning,shapes.geometric,calc,fit,mindmap,trees,backgrounds}
\usepackage{pgfplots}
\pgfplotsset{compat=1.18}
\usepackage[most]{tcolorbox}
\usepackage{listings}
\usepackage{fancyhdr}
\usepackage{titlesec}
\usepackage{hyperref}
\usepackage{cleveref}

% Paleta da casa
\definecolor{azulPetroleo}{HTML}{0B3C5D}
\definecolor{azulMedio}{HTML}{156082}
\definecolor{azulClaro}{HTML}{E8F3F8}
\definecolor{ambar}{HTML}{D99A2B}
\definecolor{ambarClaro}{HTML}{FFF4D6}
\definecolor{cinzaTexto}{HTML}{28323C}
\definecolor{cinzaClaro}{HTML}{F4F6F8}
\definecolor{verdeAcerto}{HTML}{247A52}
\definecolor{vermelhoErro}{HTML}{A13D3D}

\hypersetup{
  colorlinks=true,
  linkcolor=azulPetroleo,
  urlcolor=azulMedio,
  citecolor=azulPetroleo,
  pdftitle={TCE-MA 2026 — Auditor de Controle Externo — Tecnologia da Informação},
  pdfauthor={Material de estudo},
  pdfsubject={Concurso TCE-MA 2026 — Auditor/TI}
}

\setlength{\parindent}{0pt}
\setlength{\parskip}{0.45em}
\setlist[itemize]{leftmargin=1.5em,itemsep=0.2em,topsep=0.3em}
\setlist[enumerate]{leftmargin=1.7em,itemsep=0.25em,topsep=0.3em}

\titleformat{\chapter}[display]
  {\bfseries\color{azulPetroleo}}
  {\filleft\Large\chaptername\ \thechapter}
  {0.4em}
  {\titlerule[1.2pt]\vspace{0.5em}\Huge}
\titleformat{\section}{\Large\bfseries\color{azulPetroleo}}{\thesection}{0.6em}{}
\titleformat{\subsection}{\large\bfseries\color{azulMedio}}{\thesubsection}{0.6em}{}

\pagestyle{fancy}
\fancyhf{}
\fancyhead[L]{\small\color{azulPetroleo}TCE-MA 2026 — Auditor/TI}
\fancyhead[R]{\small\color{azulPetroleo}\leftmark}
\fancyfoot[C]{\small\thepage}
\renewcommand{\headrulewidth}{0.5pt}
\renewcommand{\footrulewidth}{0.4pt}

\tcbset{
  enhanced,
  breakable,
  boxrule=0.7pt,
  arc=1.5mm,
  left=2mm,right=2mm,top=1.5mm,bottom=1.5mm,
  before skip=7pt,after skip=7pt
}

\newtcolorbox{teoria}[1][]{colback=azulClaro,colframe=azulPetroleo,title={Teoria essencial\ifx\relax#1\relax\else: #1\fi},fonttitle=\bfseries}
\newtcolorbox{exemplo}[1][]{colback=cinzaClaro,colframe=azulMedio,title={Exemplo\ifx\relax#1\relax\else: #1\fi},fonttitle=\bfseries}
% Alias para compatibilidade com trechos técnicos redigidos com o nome em inglês.
\newtcolorbox{example}[1][]{colback=cinzaClaro,colframe=azulMedio,title={Exemplo\ifx\relax#1\relax\else: #1\fi},fonttitle=\bfseries}
\newtcolorbox{cebraspe}[1][]{colback=ambarClaro,colframe=ambar,title={Cebraspe pode cobrar assim\ifx\relax#1\relax\else: #1\fi},fonttitle=\bfseries}
\newtcolorbox{atencao}[1][]{colback=white,colframe=vermelhoErro,title={Atenção\ifx\relax#1\relax\else: #1\fi},fonttitle=\bfseries}
\newtcolorbox{resumo}[1][]{colback=white,colframe=verdeAcerto,title={Revisão rápida\ifx\relax#1\relax\else: #1\fi},fonttitle=\bfseries}
\newtcolorbox{discursiva}[1][]{colback=white,colframe=azulPetroleo,title={Treino discursivo\ifx\relax#1\relax\else: #1\fi},fonttitle=\bfseries}
\newtcolorbox{mapabox}[1][]{colback=white,colframe=azulPetroleo,title={Mapa mental funcional\ifx\relax#1\relax\else: #1\fi},fonttitle=\bfseries,breakable}

% Níveis de dificuldade exibidos apenas por estrelas: N1=★ ... N4=★★★★.
% Os identificadores N1--N4 permanecem somente no código-fonte para facilitar manutenção.
\newcommand{\tceestrela}{\textcolor{ambar}{\ensuremath{\bigstar}}}
\ExplSyntaxOn
\cs_new:Npn \tce_nivel:n #1
  {
    \str_case:nnF {#1}
      {
        {N1}{\tceestrela}
        {N2}{\tceestrela\,\tceestrela}
        {N3}{\tceestrela\,\tceestrela\,\tceestrela}
        {N4}{\tceestrela\,\tceestrela\,\tceestrela\,\tceestrela}
      }
      {#1}
  }

% Prepara o título da questão sem exibir a expressão "Certo ou Errado".
\bool_new:N \l_tce_questao_ce_bool
\tl_new:N \l_tce_questao_titulo_tl
\cs_new_protected:Npn \tce_prepara_titulo_questao:n #1
  {
    \tl_set:Nn \l_tce_questao_titulo_tl {#1}
    \bool_set_false:N \l_tce_questao_ce_bool
    \tl_if_in:NnT \l_tce_questao_titulo_tl {Certo~ou~Errado}
      {\bool_set_true:N \l_tce_questao_ce_bool}
    \tl_if_in:NnT \l_tce_questao_titulo_tl {Certo/Errado}
      {\bool_set_true:N \l_tce_questao_ce_bool}
    \tl_replace_all:Nnn \l_tce_questao_titulo_tl {Cebraspe/Certo~ou~Errado} {Cebraspe}
    \tl_replace_all:Nnn \l_tce_questao_titulo_tl {CEBRASPE/Certo~ou~Errado} {CEBRASPE}
    \tl_replace_all:Nnn \l_tce_questao_titulo_tl {Cebraspe/Certo/Errado} {Cebraspe}
    \tl_replace_all:Nnn \l_tce_questao_titulo_tl {CEBRASPE/Certo/Errado} {CEBRASPE}
    \tl_replace_all:Nnn \l_tce_questao_titulo_tl {Certo~ou~Errado} {}
    \tl_replace_all:Nnn \l_tce_questao_titulo_tl {Certo/Errado} {}
  }
\cs_new:Npn \tce_titulo_questao: {\tl_use:N \l_tce_questao_titulo_tl}
\cs_new:Npn \tce_opcoes_ce:
  {
    \bool_if:NT \l_tce_questao_ce_bool
      {
        \par\medskip
        \noindent(\hspace{1em})~certo\par
        \noindent(\hspace{1em})~errado
      }
  }
\ExplSyntaxOff

\newcommand{\nivel}[1]{\csname tce_nivel:n\endcsname{#1}}
\newcommand{\termo}[1]{\textbf{\color{azulPetroleo}#1}}
\newcommand{\certo}{\textcolor{verdeAcerto}{\textbf{CORRETA}}}
\newcommand{\errado}{\textcolor{vermelhoErro}{\textbf{INCORRETA}}}

% Questões usam título obrigatório: \begin{questao}{...}.
% Itens de julgamento recebem automaticamente as duas opções ao final do enunciado.
\newenvironment{questao}[1]{%
  \csname tce_prepara_titulo_questao:n\endcsname{#1}%
  \begin{tcolorbox}[
    colback=white,
    colframe=azulMedio,
    title={Questão \csname tce_titulo_questao:\endcsname},
    fonttitle=\bfseries
  ]
}{%
  \csname tce_opcoes_ce:\endcsname
  \end{tcolorbox}
}

% Resoluções também usam identificador obrigatório no título.
\newenvironment{solucao}[1]{%
  \begin{tcolorbox}[
    colback=cinzaClaro,
    colframe=verdeAcerto,
    title={Resolução #1},
    fonttitle=\bfseries
  ]
}{%
  \end{tcolorbox}
}

\lstset{
  basicstyle=\ttfamily\small,
  breaklines=true,
  frame=single,
  rulecolor=\color{azulMedio},
  backgroundcolor=\color{cinzaClaro},
  keywordstyle=\bfseries\color{azulPetroleo},
  commentstyle=\itshape,
  showstringspaces=false,
  columns=fullflexible
}

% Compatibilidade com mapas legados; os capítulos novos usam \mapafuncional.
\newcommand{\mapamental}[2]{%
\begin{mapabox}[#1]
\begin{center}
\begin{tikzpicture}[mindmap,concept color=azulPetroleo,text=white,level 1 concept/.append style={sibling angle=60,font=\small},every node/.style={align=center}]
  \node[concept] {#1} #2;
\end{tikzpicture}
\end{center}
\end{mapabox}
}

% Mapa funcional: estrutura + relações de prova + pegadinhas + checklist.
% #1 título; #2 núcleo; #3 estrutura/hierarquia; #4 decisões/relações; #5 pegadinhas; #6 checklist.
\newcommand{\mapafuncional}[6]{%
\begin{mapabox}[#1]
\begin{center}
\begin{tikzpicture}[
  core/.style={draw=azulPetroleo,fill=azulPetroleo,text=white,rounded corners,minimum width=9.3cm,minimum height=1cm,align=center,font=\bfseries},
  branch/.style={draw=azulMedio,fill=azulClaro,rounded corners,text width=4.5cm,minimum height=1.6cm,align=left,font=\small},
  trap/.style={draw=ambar,fill=ambarClaro,rounded corners,text width=4.5cm,minimum height=1.6cm,align=left,font=\small}
]
\node[core] (n) {#2};
\node[branch,below left=8mm and -2mm of n.south] (e) {\textbf{Estrutura / hierarquia}\\#3};
\node[branch,below right=8mm and -2mm of n.south] (d) {\textbf{Relações e decisões de prova}\\#4};
\node[trap,below=7mm of e] (p) {\textbf{Pegadinhas recorrentes}\\#5};
\node[branch,below=7mm of d] (c) {\textbf{Checklist de revisão ativa}\\#6};
\draw[-{Latex[length=2mm]},thick,azulPetroleo] (n.south) -- ++(0,-3mm) -| (e.north);
\draw[-{Latex[length=2mm]},thick,azulPetroleo] (n.south) -- ++(0,-3mm) -| (d.north);
\draw[-{Latex[length=2mm]},thick,azulMedio] (e.south) -- (p.north);
\draw[-{Latex[length=2mm]},thick,azulMedio] (d.south) -- (c.north);
\end{tikzpicture}
\end{center}
\end{mapabox}
}